\section{Empirical Materials}
\label{app:empirical-materials}

This appendix fixes the experimental objects that are part of the empirical
protocol but would interrupt the main narrative: structured resource grids, exact task
formulas, maximal-cache reuse, paired randomization, baseline matching, and
reported diagnostics.  No numerical result is inferred from this appendix;
the final release will additionally contain the realized random seeds,
measurement caches, fit outputs, and scripts producing every figure and table.
The source project also carries a theorem-to-code implementation contract that
freezes indexing, initialization, measurement, readout, selection, and paired
comparison semantics before numerical results are generated.

\subsection{Resource grids and fixed reference settings}
\label{app:dyadic-grids}

All one-dimensional resource sweeps use powers of two (with zero added only where it is a meaningful boundary value), except for qubit width.  Width is swept densely over the small hardware-relevant interval $3\leq n\leq12$ so that the empirical onset of projection-information loss is not hidden between widely spaced dyadic widths.  The planned grids are
\begin{align*}
\mathcal G_L
    &=\{0,1,2,4,8,16,32,64\},\\
\mathcal G_{\tau_+}
    &=\{2,4,8,16,32,64\},\\
\mathcal G_R
    &=\{1,2,4,8,16,32,64\},\\
\mathcal G_H
    &=\{1,2,4,8\},\\
\mathcal G_d
    &=\{1,2,4,8,16,32,64\},\\
\mathcal G_n
    &=\{3,4,5,6,7,8,9,10,11,12\},\\
\mathcal G_p
    &=\{1,2,4,8,16,32\},\\
\mathcal G_\iota
    &=\{1,2,4,8,16\},\\
\mathcal G_M
    &=\{1,2,4,8,16,32,64,128,256,512,1024,2048,4096,8192\},\\
\mathcal G_w
    &=\{1,2,4,8,16,32,64,128,256,512\},\\
\mathcal G_g
    &=\{0,1,2,4,8,16,32,64\},\\
\mathcal G_N
    &=\{128,256,512,1024,2048,4096\},\\
\mathcal G_S
    &=\{1,2,4,8,16,32,64,128,256\},\\
\mathcal G_{\tau_{\mathrm{dep}}}
    &=\{1,2,4,8,16\}.
\end{align*}
The independent input process $\rho_{\mathrm{in}}=0$ is included separately
from $\mathcal G_{\tau_{\mathrm{dep}}}$.  For a positive dependence half-life,
$\rho_{\mathrm{in}}=2^{-1/\tau_{\mathrm{dep}}}$.  The memory endpoint is
$\lambda_+=e^{-1/\tau_+}$, so every value in
$\mathcal G_{\tau_+}$ lies strictly above the lower endpoint
$\lambda_0=e^{-1}$.

Unless a campaign varies the corresponding object, the reference settings are
$\arch_{\mathrm{ref}}=(8,16,e^{-1/32})$, $N_0=1024$,
$w_{\mathrm{ref}}=512$, $g_{\mathrm{ref}}=0$, causal-proxy burn-in $B=1024$,
$S=1$, exact quantum features, and an i.i.d. input drive.  At the largest planned memory
endpoint $\tau_+=64$, the worst-case branch-state discrepancy between this
burned-in proxy and the infinite-past causal state is at most
$2e^{-16}\approx2.25\times10^{-7}$; the corresponding risk residual is the
quantity $\epsBurn(\arch,B)$ defined in
\cref{eq:burnin-risk-residual}.  Frozen-pool selection uses the
full-sample ERM rule in \Eqref{eq:frozen-pool-erm}.  Kernel smoothness, length
scale, and the Ivanov radius
are fixed once before the architecture sweeps and are not retuned separately
for each feature dimension; every method uses the same predeclared RMS-distance
rule $\norm{z-z'}_2/\sqrt p$ for its own feature dimension $p$.

\subsection{The 36-task Legendre atlas}
\label{app:task-atlas}

Every atlas target is deterministic and bounded by one, so
$Y_t=F^\star(X_{-\infty:t})$ and $\risk^\star=0$.  Coefficients are normalized
by the sum of their absolute values as in
\cref{eq:empirical-target-functional}.  The atlas contains six families.

\paragraph{Memory-depth family: eight tasks.}
For every $L\in\mathcal G_L$,
\begin{equation}
    F^{\mathrm{mem}}_L(\vx^-):=P_2(x_{-L,1}).
    \label{eq:atlas-memory}
\end{equation}
Only the oldest relevant lag changes.

\paragraph{Multiscale-mode family: four tasks.}
For $H\in\mathcal G_H$, define
\begin{equation}
    F^{\mathrm{multi}}_H(\vx^-)
       :=\frac1H\sum_{\ell\in\mathcal L_H}P_2(x_{-\ell,1}),
    \label{eq:atlas-multiscale}
\end{equation}
with
\[
\mathcal L_1=\{8\},\qquad
\mathcal L_2=\{1,64\},\qquad
\mathcal L_4=\{1,4,16,64\},\qquad
\mathcal L_8=\{0,1,2,4,8,16,32,64\}.
\]
Thus increasing $H$ increases the number of simultaneously required temporal
modes while keeping degree two and one active coordinate per mode.

\paragraph{Spatial-load family: seven tasks.}
For $d\in\mathcal G_d$,
\begin{equation}
    F^{\mathrm{sp}}_d(\vx^-)
      :=\frac1d\sum_{j=1}^dP_2(x_{0,j}).
    \label{eq:atlas-spatial}
\end{equation}
Every ambient coordinate is label-relevant, so increasing $d$ cannot be made
trivial by placing the target in one fixed low-dimensional subspace.

\paragraph{Nonlinear-degree family: six tasks.}
For $p\in\mathcal G_p$,
\begin{equation}
    F^{\mathrm{deg}}_p(\vx^-):=P_p(x_{-4,1}).
    \label{eq:atlas-degree}
\end{equation}
The active lag and arity remain fixed while polynomial order changes.

\paragraph{Interaction-arity family: five tasks.}
Fix ambient dimension $d=16$ and lag four.  For
$\iota\in\mathcal G_\iota$, define
\begin{equation}
    F^{\mathrm{int}}_\iota(\vx^-)
       :=\prod_{j=1}^{\iota}
         P_{16/\iota}(x_{-4,j}).
    \label{eq:atlas-interaction}
\end{equation}
Because every $\iota$ divides $16$, all five targets have total polynomial
degree sixteen.  The ambient dimension and memory lag are fixed; what changes
is how that degree is distributed over distinct active coordinate factors.

\paragraph{Composite family: six tasks.}
The final cases combine difficulty axes rather than changing only one:
\begin{align}
F^{\mathrm{C1}}(\vx^-)
   &:=P_8(x_{-64,1}),
   &&\text{long-memory nonlinear},
   \label{eq:atlas-c1}\\
F^{\mathrm{C2}}(\vx^-)
   &:=P_{16}(x_{-32,1}),
   &&\text{long-memory high-order},
   \label{eq:atlas-c2}\\
F^{\mathrm{C3}}(\vx^-)
   &:=\frac1{16}\sum_{j=1}^{16}P_4(x_{0,j}),
   &&\text{current high-spatial nonlinear},
   \label{eq:atlas-c3}\\
F^{\mathrm{C4}}(\vx^-)
   &:=\frac1{16}\sum_{j=1}^{16}P_2(x_{-32,j}),
   &&\text{long-memory high-spatial},
   \label{eq:atlas-c4}\\
F^{\mathrm{C5}}(\vx^-)
   &:=\frac18\sum_{q=1}^{8}P_2(x_{-\ell_q,q}),
   \quad (\ell_q)_{q=1}^8=(0,1,2,4,8,16,32,64),
   &&\text{multiscale spatial},
   \label{eq:atlas-c5}\\
F^{\mathrm{C6}}(\vx^-)
   &:=\frac14\sum_{q=1}^{4}
       P_2(x_{-\ell_q,2q-1})P_2(x_{-\ell_q,2q}),
   \quad (\ell_q)_{q=1}^4=(1,4,16,64),
   &&\text{multiscale interactions}.
   \label{eq:atlas-c6}
\end{align}
Together the six families contain
$8+4+7+6+5+6=36$ targets.  The one-axis families identify mechanisms; the
composites test whether the same resource ordering survives when several
sources of difficulty coexist.

\subsection{Mixer-mechanism controls}
\label{app:mixer-controls}

The mixer ablation in Campaign~I is deliberately narrower than a generic
circuit search.  For every paired random seed we evaluate exactly three
versions of the same branch: the identity mixer $W_n=I$, the sampled
single-qubit layer only $W_n=V$, and the full cycle mixer
$W_n=E_{\tau(J_n)}\cdots E_{\tau(1)}V$.  The latter two reuse the same sampled
single-qubit axes and angles; the full mixer adds the sampled edge gates.  No
variant receives a separately tuned readout or memory-rate budget.

Each mixer is evaluated with both the vertex-only bank $\mathcal O_V(n)$ and
the complete bank $\mathcal O_n$.  We use the fixed mechanism architecture
$n=8$, $R=8$, and $\tau_+=8$; no claim depends on this being an optimal
architecture.  For each paired repetition, first draw and freeze the shared
$\projection_n$, write
$\alpha_j(x)=\vartheta((\projection_nx)_j)$, and then generate the three
bounded current-time targets
\begin{align}
    F_{1}^{\mathrm{mix}}(\vx^-)
      &:=\cos\alpha_1(x_0),
      \label{eq:app-mixer-task1}\\
    F_{2}^{\mathrm{mix}}(\vx^-)
      &:=\sin\alpha_1(x_0)\sin\alpha_2(x_0),
      \label{eq:app-mixer-task2}\\
    F_{3}^{\mathrm{mix}}(\vx^-)
      &:=\cos\alpha_1(x_0)\sin\alpha_2(x_0)\sin\alpha_3(x_0).
      \label{eq:app-mixer-task3}
\end{align}
All six mixer/bank variants share that preprocessor and label sequence.  The
first target is a one-site control; the second is directly compatible with an
identity edge observable and with the mixed vertex term in
\cref{cor:mixer-direct-interactions}; the third matches the mixed edge-seed
three-stream term.  Defining these targets after the independent preprocessor
draw is intentional: they test the theorem in its natural projected
coordinates without letting Gaussian preprocessing itself confound the
interaction-order ablation.  They are therefore reported only as mechanism
diagnostics and are kept separate from the preprocessor-independent 36-task
atlas.

To test whether the same mechanism matters beyond targets tailored to the
encoding coordinates, we repeat the identical six mixer/bank variants on the
three fixed atlas tasks
\[
    F^{\mathrm{mem}}_8,\qquad
    F^{\mathrm{int}}_{16},\qquad
    F^{\mathrm{C6}}.
\]
These labels are generated independently of the Gaussian preprocessor.  The
first is a low-interaction negative control, while the latter two require
higher-order spatial or spatiotemporal interactions.  No mixer-specific
hyperparameter is changed between the theorem-aligned and atlas-aligned
slices.

The complete and vertex-only observable banks also receive one finite-shot
cost-matched comparison using the full cycle mixer on $F^{\mathrm{C6}}$.  A
dyadic value $M_{\mathrm E}^{\mathrm{cmp}}$ is
frozen in the pre-run manifest.  The complete bank uses all nine global Pauli
settings at $M_{\mathrm E}^{\mathrm{cmp}}$ repetitions per setting.  The
vertex-only bank requires only the three all-$X$, all-$Y$, and all-$Z$ settings
and uses $M_{\mathrm V}^{\mathrm{cmp}}=3M_{\mathrm E}^{\mathrm{cmp}}$
repetitions per setting.  Thus both variants consume exactly
$9M_{\mathrm E}^{\mathrm{cmp}}$ branch-circuit runs per input window.  Training
and held-out measurement caches remain independent exactly as in Campaign~II.

For exact circuit diagnostics, expand
$W_n^\dagger O W_n=\sum_P c_{O,P}P$ in the normalized Pauli basis and report
\begin{equation}
    \spreadScore(W_n;\mathcal O)
      :=\frac{1}{|\mathcal O|}
        \sum_{O\in\mathcal O}
        \sum_{\operatorname{supp}(P)\not\subseteq\operatorname{supp}(O)}
          |c_{O,P}|^2.
    \label{eq:app-pauli-spread-score}
\end{equation}
The score is identically zero for $W_n=I$ and for the local-only $V$ layer,
whereas the graph-local spreading result predicts positive mass outside the
seed support for the relevant full-mixer observables almost surely.  We also record the stable rank of the column-centered exact feature
matrix on a common input sample, using exactly
$r_{\mathrm{stab}}(F_c)=\norm{F_c}_{\mathrm F}^2/\norm{F_c}_2^2$ with
$r_{\mathrm{stab}}(0)=0$, as defined in \Eqref{eq:stable-feature-rank}.  These diagnostics test the proposed
mechanism---operator spreading and observable enrichment---without treating
performance alone as proof that entanglement is useful.

\subsection{Maximal-cache construction and paired randomness}
\label{app:maximal-cache}

The large grids are made economical by coupling conditions without changing
any required marginal distribution.  The maximal caches are campaign-local:
$R_{\max}$ for the architecture study, $M_{\max}$ for the $S=1$ measurement
study, and $(S_{\max},R_{\mathrm{search},\max})$ for the exact-feature search
study.  They are never multiplied into one unnecessary maximal Cartesian
cache.

\paragraph{Reservoir prefixes.}
For every Campaign~I design, draw $R_{\max}=64$ independent branch seeds
$(U_r,\omega_r)$ with $U_r\sim\operatorname{Unif}(0,1)$.  At memory endpoint
$\lambda_+$ set
\begin{equation}
    \lambda_r(\lambda_+)
      :=\lambda_0
        \left(\frac{\lambda_+}{\lambda_0}\right)^{U_r}.
    \label{eq:coupled-lambda-draw}
\end{equation}
For each fixed $\lambda_+$ this has exactly the density in
\cref{eq:log-uniform-memory-law}, while reuse of the same $U_r$ gives a paired
comparison across memory endpoints.  The architecture with $R$ branches uses
only the first $R$ pairs.  Consequently all values in $\mathcal G_R$ are
available from one $R_{\max}$ draw.

\paragraph{Frozen-pool prefixes.}
Generate $S_{\max}=256$ candidate designs once for Campaign~III, each with
$R_{\mathrm{search},\max}=32$ exact-feature branches.  The setting with pool
size $S$ and branch count $R$ uses candidates $1{:}S$ and branches $1{:}R$,
then applies the identical full-sample ERM selection rule.  Oracle-best curves are computed
on the same prefixes.  This preserves the order-statistic nesting of random
search, yields every
$S\in\{1,2,4,8,16,32,64,128,256\}$ and
$R\in\{1,2,4,8,16,32\}$ from one maximal exact-feature pool, and removes
irrelevant between-pool Monte Carlo noise.

The compact selection-pressure slice uses the same $S_{\max}=256$ candidate
ordering at fixed $(n,R,\tau_+)=(8,16,32)$, with
$N\in\{256,1024,4096\}$ and $S\in\{1,16,64,256\}$.  Within a paired repetition,
the $N=256$ and $N=1024$ training sets are prefixes of the $N=4096$ gapped
training realization, so increasing $N$ changes only the amount of observed
data.  Held-out evaluation uses one common independent test realization across
all $(N,S)$ cells.  The slice is run only on
$F^{\mathrm{multi}}_8$ and $F^{\mathrm{C6}}$ and uses exact features, so it
isolates finite-pool selection from measurement noise.

\paragraph{Measurement prefixes.}
For Campaign~II, where $S=1$, generate $M_{\max}=8192$ independent outcomes
for every input window, tested branch, and one of the nine global Pauli
settings.  The feature estimate at repetition count $M$ uses outcomes $1{:}M$.
Since a single global outcome supplies all compatible Pauli coordinates, this
cache preserves the within-shot correlations of the grouped measurement
scheme.  The total acquisition required to support an entire $M$ sweep at a
fixed tested architecture is therefore the maximal cache once, not
$\sum_{M\in\mathcal G_M}9NMR$ separate acquisitions.  The representative
finite-shot pool-ranking check in Campaign~III is generated separately only up
to $S=64$ at its chosen dyadic $M$; it does not instantiate the full
$S_{\max}M_{\max}$ product.

\paragraph{Coupled Gaussian preprocessors.}
For a task of ambient dimension $d$, draw one matrix
$G\in\sR^{n_{\max}\times d}$ with $n_{\max}=12$ and i.i.d. standard-normal
entries.  For $n\in\mathcal G_n$, set
\begin{equation}
    \projection_n:=\frac1{\sqrt n}G_{1:n,:}.
    \label{eq:coupled-preprocessor}
\end{equation}
Each $\projection_n$ has the exact marginal law in \cref{eq:gaussian-projection},
while widths share as much preprocessing randomness as possible.  The random
mixing circuit itself is still sampled from the correct width-specific law
$\mathbb P_n^\omega$; width comparisons are paired by task, temporal sample,
preprocessor seed, and random-seed stream rather than interpreted as nested
quantum circuits.

\paragraph{Finite-window causal-proxy construction.}
For the window-fidelity experiment, the same temporal trajectory is propagated
through each branch after a burn-in $B=1024$.  The resulting state is a
\emph{certified causal proxy}, not an exact infinite-past state.  Contraction
bounds its branchwise trace-distance residual by $2\lambda_+^B$, and the same
feature/readout argument as in the finite-window theorem bounds the associated
risk discrepancy by $\epsBurn(\arch,B)$.  At the largest planned
$\tau_+=64$ this branch-state residual is $2e^{-16}\approx2.25\times10^{-7}$.
The finite-$w$ state is obtained by resetting to $\rho_+$ exactly $w$ steps
before the same label time and replaying the identical suffix.  Every reported
window gap is accompanied by $\epsBurn(\arch,B)$ so numerical burn-in error
cannot be mistaken for finite-window error.

\subsection{Homogeneous-memory control}
\label{app:homogeneous-control}

The heterogeneous-memory comparison must preserve branch count and frozen
mixing diversity.  For a given $R$, retain the same $\omega_1,\ldots,\omega_R$
but replace the independently drawn rates by one common rate.  The primary
control uses the geometric center of the log-rate interval,
\begin{equation}
    \lambda_{\mathrm{hom}}
      :=\sqrt{\lambda_0\lambda_+},
    \label{eq:homogeneous-rate}
\end{equation}
so no extra tuning budget is introduced.  On the representative multiscale
slices used for the architectural claim, we additionally evaluate
$\tau_{\mathrm{hom}}\in\mathcal G_{\tau_+}$ and choose the common rate by the same full-sample empirical-risk rule.  A heterogeneous multiplex must therefore be
competitive not only with one fixed common memory scale but also with a
full-sample-ERM-selected homogeneous alternative.

\subsection{Baselines and matched resources}
\label{app:empirical-baselines}

\paragraph{Projection-only Mat\'ern control.}
This control is used only in the width-versus-spatial-load study.  For the same
shared Gaussian matrix used by QuaRK at width $n$, it receives
\[
    (\projection_n x_{-w_{\mathrm{ref}}+1},\ldots,\projection_n x_0)
    \in\sR^{nw_{\mathrm{ref}}},
\]
flattens the projected window in time order, and fits the same fixed
RMS-normalized Mat\'ern/Ivanov readout.  It has no reservoir dynamics, mixer,
or observable bank.  Its purpose is not to compete as a third general
baseline, but to expose empirically how much predictive information survives
$\projection_n$ before the quantum representation acts.  It therefore uses the
same projection realization, temporal sample, kernel parameters, Ivanov radius,
and held-out windows as the corresponding QuaRK width condition.

\paragraph{Direct finite-window Mat\'ern.}
The raw baseline receives the same observed history $x_{-w_{\mathrm{ref}}+1:0}$,
flattened in time and coordinate order, and uses the same normalized Mat\'ern
kernel family, the same RMS-distance rule $\norm{z-z'}_2/\sqrt p$ for its own
feature dimension $p$, and the same fixed Ivanov policy.  It tests whether the
controlled target requires a reservoir representation at all.

\paragraph{Classical multiscale reservoir.}
The classical reservoir consists of $R$ independent frozen leaky branches
driven by the same input and concatenated before the same Mat\'ern readout.
Its forgetting parameters are matched branchwise to the QuaRK $\lambda_r$
values, and its total exposed state dimension is matched to
$m_{\arch}=12Rn$ up to integer allocation across branches.  Input and recurrent
weights are frozen random draws; no internal parameter is trained.  The
same RMS-distance rule is applied to its concatenated feature vector, so
changing exposed dimension does not trigger a separate kernel retuning.  The
comparison therefore preserves the two structural ingredients most relevant to
the QuaRK claim---parallel frozen dynamics and heterogeneous memory---while
changing the quantum feature generator.

Neither baseline is intended to establish a quantum-advantage theorem.  The
raw kernel can reveal that a reservoir is unnecessary, and the classical
multiscale reservoir can reveal that any observed gain follows from generic
multiscale reservoir structure rather than specifically quantum features.

\subsection{Diagnostics and paired reporting}
\label{app:empirical-diagnostics}

Every resource curve uses paired temporal samples and paired design seeds
where the construction above permits it.  Except for the training-risk side of
$\widehat\Delta_{\mathrm{gen}}$, all reported predictive risks are evaluated on
held-out windows; paired quantities use the same held-out windows and labels,
and any measured evaluation uses measurement randomness independent of the
training cache.  The final numerical study should use
at least sixteen paired repetitions for architecture and fidelity curves and
thirty-two for the $S$-selection curves, where order-statistic variability is
larger.  A long independent test realization is used for all population-risk
approximations.

The primary metric is \cref{eq:nemse}.  Campaign~I's width study reports the
raw-history Mat\'ern baseline, projection-only Mat\'ern control, and QuaRK on
the same paired width conditions.  Its mixer mechanism slice additionally
reports $\spreadScore$ from \cref{eq:app-pauli-spread-score} and exact-feature
stable rank for the identity, local-only, and full-cycle mixers under
vertex-only and complete observable banks.  The same mixer/bank comparison is
reported separately on the three preprocessor-independent atlas controls, and
the finite-shot vertex/full-bank check is compared at identical branch-circuit
run count.  Campaign~II reports
$\widehat\Delta_{\mathrm{shot}}$,
$\widehat\Delta_{\mathrm{win}}$, and
$\widehat\Delta_{\mathrm{gen}}$, always evaluating the same fitted readout on
the two feature/risk objects being compared.  Campaign~III reports ERM-selected NEMSE, oracle-best-in-prefix NEMSE,
selection regret, and $C_Q=9NMSR$.  Its compact $N\times S$ slice reports the
same oracle and regret quantities as a surface or grouped curves, so the effect
of additional data on ERM over-selection is visible directly.  The observed
regret is the empirical check on the finite-sample selection price that is
controlled theoretically by the uniform $\log S$ term; the certificate terms
are not tabulated merely because they exist.

Finally, because the feature dimension $12Rn$ changes with architecture, the
release records the empirical distribution of pairwise Mat\'ern similarities
computed with the predeclared RMS metric at each $(n,R)$ setting.  This is a
sanity diagnostic for residual kernel saturation, not another tuned
hyperparameter.  The kernel specification is not changed after inspecting the
diagnostic: any residual collapse toward uniformly zero or uniformly one
similarity is reported as a limitation of that fixed readout specification,
not repaired by architecture-specific length-scale tuning.

The final empirical release will contain the exact generator implementation,
all task identifiers from \cref{app:task-atlas,app:mixer-controls}, temporal and quantum random
seeds, maximal branch/pool/measurement caches, baseline implementations,
resource configurations, train/test indices, and the scripts that map
those objects to every reported figure and table.
